%; whizzy chapter
% -initex iniptex -latex platex -format platex -bibtex jbibtex -fmt fmt

%   Pdf作成手順
% dvipdfmx debianmeetingresume200506.dvi


%     Tokyo Debian Meeting resources
%     Copyright (C) 2006 Junichi Uekawa

%     This program is free software; you can redistribute it and/or modify
%     it under the terms of the GNU General Public License as published by
%     the Free Software Foundation; either version 2 of the License, or
%     (at your option) any later version.

%     This program is distributed in the hope that it will be useful,
%     but WITHOUT ANY WARRANTY; without even the implied warranty of
%     MERCHANTABILITY or FITNESS FOR A PARTICULAR PURPOSE.  See the
%     GNU General Public License for more details.

%     You should have received a copy of the GNU General Public License
%     along with this program; if not, write to the Free Software
%     Foundation, Inc., 51 Franklin St, Fifth Floor, Boston, MA  02110-1301 USA

%この資料はA3用紙に両面印刷で印刷する予定

\documentclass[mingoth]{jsarticle}
\usepackage[dvipdfmx]{graphicx}
\usepackage{fancybox}
\usepackage{longtable}
\usepackage[dvipdfmx]{hyperref}
\usepackage{url}

\AtBeginDvi{\special{pdf:tounicode EUC-UCS2}}

% とりあえずcommandline環境を定義。入出力についてはcommandline環境を活用
%する
\newenvironment{commandline}%
{\VerbatimEnvironment
  \begin{Sbox}\begin{minipage}{12cm}\begin{tiny} \begin{BVerbatim}}%
{\end{BVerbatim}\end{tiny}\end{minipage}\end{Sbox}
  \setlength{\fboxsep}{8pt}\fbox{\TheSbox}}



% 三択問題用
\newcounter{santakucounter}
\newcommand{\santaku}[4]{%
\addtocounter{santakucounter}{1}

\nopagebreak 問題\arabic{santakucounter}. 
#1\\
\nopagebreak□ A #2\\
\nopagebreak□ B #3\\
\nopagebreak□ C #4
\pagebreak[1]
\hspace{1cm}
\\

}

\newcommand{\emptyspace}{(\underline{\hspace{1cm}})}


% sectionをセンタリングする
\makeatletter
  \renewcommand{\section}{\@startsection{section}{1}{\z@}%
    {\Cvs \@plus.5\Cdp \@minus.2\Cdp}% 前アキ
    {.5\Cvs \@plus.3\Cdp}% 後アキ
    {\normalfont\Large\headfont\raggedright\centering}} % style
\makeatother

% section の代わりの環境
\newcommand{\dancersection}[2]{%
\newpage
東京エリアDebian勉強会 2005
\hrule
\vspace{0.5mm}
\hrule
\hfill{}\includegraphics[width=3cm]{image200502/openlogo-nd.eps}\\
\vspace{-4cm}
\begin{center}
  \section{#1}
\end{center}
\hfill{}#2\hspace{3cm}\space\\
\hrule
\hrule
\vspace{1cm}
}






\begin{document}

\begin{titlepage}

% 毎月変更する部分
\title{
 第5回 東京エリア Debian 勉強会\\事前資料
\footnote{機密レベル public: 一般開示可能}}
\date{2005年6月11日}
\author{Debian勉強会会場係 上川 純一\thanks{Debian Project Official Developer}} 
\maketitle

\end{titlepage}

\newpage
\tableofcontents

\dancersection{Introduction To Debian 勉強会}{上川 純一}

今月のDebian勉強会へようこそ。
これからDebianのあやしい世界に入るという方も、すでにどっぷりとつかってい
るという方も、月に一回Debianについて語りませんか？

目的として下記の二つを考えています。

\begin{itemize}
 \item メールではよみとれない、もしくはよみとってられないような情報を情
       報共有する場をつくる
 \item まとまっていないDebianを利用する際の情報をまとめて、ある程度の塊と
       して出してみる
\end{itemize}

また、東京にはLinuxの勉強会はたくさんありますので、Debianに限定した勉強
会にします。Linuxの基本的な利用方法などが知りたい方は、他でがんばってくださ
い。
Debianの勉強会ということで究極的には参加者全員がDebian Packageを
がりがりと作りながらスーパーハッカーになれるような姿を妄想しています。

Debianをこれからどうするという能動的な展開への土台としての空間を提供し、
情報の共有をしたい、というのが目的です。
次回は違うこと言ってるかもしれませんが、御容赦を。

\subsection{講師紹介}

\begin{itemize}
 \item{kmutoさん} debian-installer ハッカーです。
 \item{えとーさん} dsysの開発者です。
 \item{上川 純一} 宴会の幹事です。
      Debian Developerです。
      元超並列計算機やっていて、
      今は音楽関係とか、気づいたらcannaとか。
      あと、pbuilderや、libpkg-guideなどを通して、Debianの品質向上を目指してます。
\end{itemize}

\subsection{事前課題紹介}

今回の事前課題は
「debian-installerをつかってみて/今後のdebian-installerに期待するもの」
というタイトルで200-800文字程度の文章を書いてください。
というものでした。
その課題に対して下記の内容を提出いただきました。

\subsubsection{みつかさん}

少し前に使ったきりなので記憶が曖昧なのですが…。
まず驚いたのが国際化が進んだなということです。
各国語のメニューがずらっと並んだなかに日本語も含まれているというのは
なんだかとても誇らしいような気持ちになりました。
日本語を選択して進んでいくと普通にすんなりとインストールできてしまい
感動しました。
以前のインストーラーと比べて、ステップ数も少なくなったと思います。
これも、インストールが簡単になったなという印象に寄与していると思います。

今後への期待としては、日本語を選択してインストールした場合

\begin{itemize}
 \item jfbtermがデフォルトで起動したり
 \item ja\underline{ }JP.eucJP辺りがデフォルトで設定されてたり
\end{itemize}


すると嬉しいかなあとか思いました。


\subsubsection{kaepapa}

debianのインストーラはpotatoとwoodyそしてsargeで使ってきたが徐々に
インストールが容易なものになってきたと思います。
個人的にはwoodyの時のインストーラで十分満足できる状態でした。
sargeになりcfdiskでなくなりちょっと残念です(経緯は知っている)。
partitionerに関してはちょっと間違うとHDD全域をパーティショニングしてしまう
リスクを感じました。
この点は十分な情報の提供が必要だと思われます。RC3当時のpartitionerの
不具合(パーティションテーブル不全？)もなくなり(20050419)
リリースされたらもう一度ためしてみたいと考えています(VMWareゲスト)。
またパッケージの選択等については以前よりインストーラでは最小構成のみ
インストールし後からaptitudeで必要なものを導入するという癖がついているため
あまり記憶に残っていませんが以前とかわらず最小構成でインストールが終れることに好感が持てます。
このポリシーはずっと継続して欲しいと願っています。

\subsubsection{やまねさん}

使ってみた感想は昔に web に載せたことがあるので略。その頃からするとぐっとよくなってます :-)

\begin{itemize}
 \item 改善して欲しいーという点。
       \begin{itemize}
	\item やはり「X の自動認識」の弱さが気になる。初心者には鬼門。
	\item パーティションの分割し辛さと RAID 機能の弱さ
	\item パッケージダウンロードをしている間に他の設定をさせるなどは難しいのか? ずーっと待ってるのはちょっと。
	\item 特定機種での問題（Dell のとか）は結局解決したんだろうか。
	\item 「定期的なアップデート」。たとえ stable といえど、新しい機種に対応していないインストーラーを勧められるのは「？」というところ。ここはインストーラだけでも半年に1回ぐらい更新を行うようなポリシーに切り替えて欲しい。
	\item iptables でパケットをデフォルト遮断は難しい? WindowsXP SP2 なども基本は『保護』の方向だし、この機能は欲しい
	\item tasksel の分け方とか、本当に「最小インストール」するのが簡単になって欲しいところとか
       \end{itemize}
 \item 思いつき
       \begin{itemize}
	\item
イベントのときに「d-i fest」でひたすらインストールできた・できない機種の情報収集するとか。
	     \begin{itemize}
	      \item
		   BTS 報告だと英語がネックで厳しいし
	      \item こんなイベントなら参加の敷居も低いし
	      \item レアマシン持ってくる人もいて、結構なデータ取れるかもしれないし
	      \item エスパー能力なくてもその場で解決できること多そうだし
	     \end{itemize}
       \end{itemize}
\end{itemize}
　ということで如何でしょうか :-)


\subsubsection{中島　清貴さん}

 Debian辞典という本がちかくの本屋さんに平積みで売っていたので、さらっと立ち読
みしてみたらインストール方法から何からすべて書いてあるではないか。これだ、この
本さえあれば今すぐDebianが使えるはずだ。さっそく買って家で読んでみる。まったく
Debianなるものとはかけはなれた業界にいるような私にも分かるような事細かな説明。
projectの組織構成まで書いてある。これは良書に間違いない。 そこで友人Nを電話で
モスバーガによびだして彼にも買わせる。附属のCDに特制インストローラが付いている
ので、これでインストールはだいじょぶなはずだ。このような本が読めることはDebian
を知っている私としての最高の祝福だと思う。

\subsubsection{澤田さん}
Debian-Installerを使ってみて感じたこととして、woody時代に 
比べてLinuxを知らない人でもインストールが容易になったなぁ 
ということがある。大学時代にサーバを作ろうということで 
Debianをインストールするというアレゲな実験がありTAをやって 
いたわけだが、多くの人がはまっていたこととして、


\begin{itemize}
 \item  パーティションの設定
 \item  ドライバの組み込み
 \item  パスワードを打っても何も表示されないこと
\end{itemize}

ってのがあった。Debian-Intallerではここら辺が改善されてい 
るのでいい感じだ。

また、はじめに日本語を選択するとキーボードとかタイムゾーンとか 
パッケージとかを日本向けにしてくれるのもとってもフレンドリーに感 
じた。

逆にちょっとあれかな〜と思ったのはネットワーク設定のデフォルトが 
DHCPな点である。ここはstaticにするかdhcpにするかの選 
択肢がほしかったなと思った。



\subsubsection{Ryoichi Kato}
第5回、Debian勉強会に参加するのにあたり,
実際に使ってみて感想を書いてみようと思い,
久々に真っさらの状態からdebian のインストールをしてみました. 
手元に潰してもいいマシンがなかったので,
CD image (sarge-i386-netinst.iso) から
VMWare 上にインストールしました.

普段あまり debian-installer を使う機会もなく,
前回 debian をインストールしたのは 2,3 年前ですが,
その頃の印象と比べると非常に使いやすくなっていると感じました.
以前あった「クセ」, 「アク」(コツを知らないと使いこなせない)
のようなものがずいぶん薄くなっていて, 次何をすれば良いんだろう」と
迷う必要が殆んどなくなったという印象です.
また, 特にパーティショニング関連は, 細い設定をするのや
やり直しをするのがかなり楽でした.

ところで, grubをインストールして再起動するところまでは
すんなり言ったのですが, その後の再起動でちょっとトラブルに見舞われました.

VMwareのディスク領域を2Gしか確保してなかったのですが,
base-config で "desktop environment" を選択しておいたら
パッケージをインストールしている最中に disk full になってしまいました.
 (最初に必要な領域が 600MB + 1700MB 位とか出たのですが,
  空き容量チェックはしてないのでしょうか?)
そのままにっちもさっちも行かなくなったので, そこで二枚目のコンソールへ移って
\begin{commandline}
#  rm /var/cache/apt/archives/*.deb!
\end{commandline}

してから手動で apt-get でインストールを続行しようとしたら,
/var/cache/debconf/config.dat が lock されてると言われてしまい
そのままでは手動でインストールできませんでした.

結局 base-config を kill -STOP (respawnなので)することで
なんとか無理矢理 install してみましたが,
このような場合はどうやるのが正解だったのでしょうか?
\subsubsection{小林儀匡さん}

以前のboot-floppiesインストーラでも質問に従っていけばインストールは問題なく
できましたが、

\begin{enumerate}
 \item 質問やメッセージなどの数がやたらと多いこと、
 \item それらがすべて英語で書かれていること、
 \item それらの中には初心者にとってよく分からないものも多いこと
\end{enumerate}
などが理由で、使い勝手はあまりよくなかったように思います。
「一応インストールしてみたものの、全てがまともに設定されたかはやや不安だ」
というような心地に、(少なくとも以前の自分は) なりました。
そのため、新しいdebian-installer (d-i) で
\begin{enumerate}
 \item  多くのプロセスが簡潔化されたこと、
 \item 質問やメッセージもすべて翻訳されたこと、
 \item 初心者がどうしていいか悩むような設定 (例えばパーティション分け) に
    いくつかの簡単な選択肢が用意されたこと
\end{enumerate}
を非常に嬉しく思います。

d-iのおかげで、サービスをほとんど入れる必要がない (taskselで何も選択しなくて
済む) 数値計算用のマシンなどは、たった10分でインストールが完了しました。
また、まだ挑戦したことはありませんが、プレシードを利用して多数のマシンに
一度にDebianを入れることが簡単にできるようなので、多数のホストからなる
計算機環境を作り上げる機会があったら試してみたいものです。

一度インストールしてしまえば再インストールを必要とせず、ハードウェアが壊れる
まで使い続けられるのを売りにするDebianなので、「Debianユーザ」になった人に
とってインストールはあまり重要ではないといっても過言ではないでしょう。
しかし一方で、そのインストールの敷居が高いために「Debianユーザ」になるのを
挫折するユーザもいると思います。
d-iによってそのような人にDebianの恩恵を与えられることを期待するとともに、
コミュニティへの新規参入者のためとも言えるようなインストーラをわざわざ作った
開発者の方々に感謝します。

自分にとってd-iは既に十分に使いやすいレベルに達しているため、今後のd-iに欲しい
新機能や改善点などは今のところ思い付きません。
あえて挙げるとすれば、WindowsやMacの世界から来た初心者の多くにはGUIのほうが
とっつきやすいと思うので、GUI化が実現できるとよいかもしれません。
もっとも、CUIに慣れた人間からするとGUIはいろいろと面倒なので、たとえGUI化
されても、CUIとGUIのどちらを使うかを最初に選べるようになっていると嬉しいです。

\subsubsection{えとーさん}

  サーバなどにインストールすることも何度かあったのですが、
その際にSoftware-RAID、LVMの構築があまり好ましくなく、いろいろと
不便でした。

  よって、dm,mdへの対応の強化が行なわれることが好ましいと
思っています。

  他にはカスタムインストーラを手軽に作れるようになると
嬉しいなぁと思っています。これは、ノートPCや、サーバ機器などの、
ハードウェア構成がある程度固定化しているものについては、一部設定を
省けるようなインストーラが欲しいからです、例えば、ハードウェアの
設定部分をテキストファイルから読み込めるようにしておき、テンプレートに
従って記述したテキストファイルをインストーラに追加しておけば、
インストーラの画面から型番等を選択することにより設定終り。

  パッケージ選択についても同様でもハードウェアによって必要なパッケージを
探すのが面倒ですので環境に合わせたパッケージのタスクみたいのを選択できる
ようになっていて、そこからも選べるようになるといいなぁと思います。

\subsubsection{岩松 信洋}
debian-installer を使ってみての感想
	Debian辞典のCDを使ってみました。


\begin{enumerate}
 \item インストール文章が全て日本語。すばらしい。
 \item ある程度自動的にパーティションを切ってくれるシステム
	  （ガイドによるパーティショニング）が手間を省いてくれてよい。
\end{enumerate}	

期待すること
\begin{enumerate}
 \item	  マウスを使えるようにして容易にインストール。（初心者向け）
 \item	  XWindowsの設定を容易にしてほしい。
 \item	  日本語でインストールしたときはFBでも日本語表示できるようにしてほしい。
	   (現状で文字化けする。)
\end{enumerate}


\subsubsection{上川}


debian-installerは実は使ったことがない。
debian-installerによってリプレースされたboot-floppiesの開発には参加して
いた。
boot-floppiesは使っていたが実際CDから起動してシェルを立ち上げるためのみに利用
して、あとはdebootstrapとかを利用し、手動でインストールしている。

さて、debian-installerはetchにむけて今後どうなっていくことが期待できるだ
ろうか。

まず、言語の選択のよりよい流れが欲しい。
日本語をインストール画面で選択したら、コンソール画面も日本語で進み、イン
ストールし終ったときのデフォルトの画面も日本語になるようにしたい。
これを実現するには 下記の二つが実現したら楽になるのではないだろうか。

\begin{itemize}
 \item \verb!ja_JP.UTF-8! localeであらゆるアプリケーションが標準で動くようになること
 \item Linuxカーネルがutf-8での日本語の表示対応したフレームバッファをドライブでき
       るようになること
\end{itemize}

etchとしてあと一年くらいハックできる期間があるのだから、これくらい実現で
きそうな気がするのだが、どうだろう？

%%% trivia quiz
\dancersection{Debian Weekly News trivia quiz}{上川 純一}

ところで、Debian Weekly News (DWN)は読んでいますか？
Debian 界隈でおきていることについて書いているDebian Weekly News.
毎回読んでいるといろいろと分かって来ますが、一人で読んでいても、解説が少
ないので、
意味がわからないところもあるかも知れません。みんなでDWNを読んでみましょう。

漫然と読むだけではおもしろくないので、DWNの記事から出題した以下の質問にこたえてみてください。
後で内容は解説します。

\subsection{2005年21号}

\santaku{12Wの電力消費で動くコンパクトなDebianサーバを準備するために、Silas
Benettはどうしたか}{となりの家の犬を無線LANステーションにした}{Mac Miniをかってきて、CDROMドライブをとりはずして
バッテリーを装着した}{Dual Xeonのサーバを購入してXeonを一つにへらしてみた}
\santaku{Michael BanckはHurdで何をしたか}{GNOMEとQTを動かした}{もうあきら
めようと宣言した}{tuxracerを動かした}
\santaku{Orphanされたパッケージなどの一覧が出て来るWNPPメールは今後どう
なるか}{毎週debian-devel-announceにメールを投げていたのを停止して、debian-wnppに
移動する}{もう orphan なんてものは存在しない}{orphanしたパッケージは今後は無視
する}
\santaku{Nico Goldehはunrarパッケージのバージョンがずいぶんさがっていることに気づきまし
た。その理由は}{もともとのunrarはnon-freeで、freeな実装ができたのでそち
らに移行した}{unrarの開発は後向きにすすんでいる}{unrarなんて知らないので
関係ない}
\santaku{Wasteパッケージにはどういうライセンス上の問題があるか}{ソフトウェ
アとしてつかいものにならない}
{開発元が一旦GPLでリリースしたが、revokeした}{ライセンスが入っていない}
\santaku{Debian woody のアップデートが出たが、そのバージョンは}{3.0r6}{3.0.6}{3.0-RC6}


\subsection{2005年22号}

\santaku{Andreas BarthがBTSのLDAPゲートウェイの高速化のために試したのは}
{Archivedバグが多すぎるので、Archivedバグを分離してしまいメモリ負荷を減
らす}{LDAPをやめてSQLにする}{不要っぽいバグはなかったことにする}
\santaku{Philipp KernはDebian Archiveにvideoセクションを追加しようと提案
しました。videoセクションが無いため現状は各ビデオ関連のパッケージは}{どこのセクショ
ンにも属していない}
{libsセクションに投げ込まれている}{該当する
アプリケーションはgraphics
セクションとデスクトップ環境のセクションに混在している}
\santaku{debian-legalのサマリーページについてFrank Lichtenheldはもう削除
しようと提案しました。その理由は}{使えないから}{面白くないから}{合意がとれないのであたらしくサマリー
が作成できずメンテナンスが不可能}

\subsection{2005年23号}

\santaku{Debianのリリースをおくらせる必要がある、と主張していたKDEのバグは何？}
{webcollageスクリーンセーバがデフォルトで動作するようになっているので、
オンラインのポルノ画像が表示される可能性がある}{
KDEが動かない}{KDEの壁紙がDebian向けではない}
\santaku{Debian GNU/Linux 3.1がリリースされたのは？}{
2005年6月6日(日本時間6月7日)}{
2005年7月1日(日本時間7月2日}{
2005年1月1日(日本時間1月2日)}
\santaku{3.1r0のCDセットの問題は}{だれも焼けないような枚数になった}{デフォ
ルトインストールのsources.listのsecurity.debian.org向けの行が有効でなかっ
た}{大きすぎて世界中のネットワーク帯域をくいつぶしてしまう}
\santaku{Wesley Landaker はGPGキーサインのための確認に実際に会わなくても
できる基準を提案しました。その反論は}{GPGって出会い系サイトでしょ？}{写真のみで確認するのであれば、写真
は偽造しやすいので、その情報を信頼するのは難しい}{会わなくてよいのだった
らGPGは面白くない}
\santaku{ライブラリのABIが変更になった場合に気を付けることは何か}{
ライブラリパッケージ名を変更しておけば管理者が明示的にアンインストールす
るまでは古いライブラリバージョンもシステムに残る}{古いライブラリにリンク
してあるプログラムを確実に全部動かなくする}{変更履歴に怒りのコメントを記
述する}

% \santaku{}{}{}{}

\dancersection{最近のDebian関連のミーティング報告}{上川 純一}

\subsection{東京エリアDebian勉強会5回目報告}


前回開催した第5回目の勉強会の報告をします。

	    DWNクイズをしました。

	    松山さんが DFSG について語りました。

	    岩松さんが dpkg-cross の使い方について説明しました。

	    上川が lintian, linda の使い方について説明しました。
	    Debianのパッケージの作成のフローを一度デモできるとよいねぇ、という話しになりました。

	    グループディスカッションでは、今後のアーキテクチャ対応をどうしましょうか、という話しをしました。
	    現在のアーキテクチャはCPU毎にわけられており、CPU以上の部分についてはパッケージでカバーしているのですが、task-ibookg4とかがあれば、各ハードウェアに必要なパッケージをごっそりいれれてよいよねぇ、という案が出たりしました。
	    全パッケージをスクリプトで再実装してしまい、
	    各アーキテクチャの差分を減らす、という迷案も出ていました。


\dancersection{alternatives -選択せよ-}{えとー}
%% えとーさんの記事はここから
\subsection{alternatives とは？}
     パッケージ管理というと依存関係の管理についてのみがクローズアップされるが、
   パッケージ管理としては、それだけでは不足する部分がある、それを補うものの一つが
   alternatives です。
  
     日本語訳すると「選択肢」、私の定義ですが、
   「複数のパッケージを特定機能ごとに一つにまとめて扱い機能ベースのパッケージの管理を提供するもの。」
   「パッケージを機能の視点で管理する。」
   「プログラムに OS の 統一的な API を提供する。」
    の3点です。

   機能としては、
     「類似の機能を持つプログラムの別名を提供する。」で、symlink を使って実現しています。
    alternatives とは、その symlink を管理するための機構と言えるでしょう。
 
   なぜわざわざ別名が必要？

\begin{enumerate}
 \item プログラムからの使い勝手の向上
     万単位のパッケージの存在する Debian では同じような機能を持つ
     別々のパッケージが複数存在することは珍しくありません。
     伝統的にUnix系OSではパイプやリダイレクトなどや
     スクリプト、プログラムなどから他のプログラムの機能を利用することにより
     コード量の削減と質の向上を図っています。
     しかし、プログラムから呼び出す際に同じような機能をもつもの全てを把握し
     条件分岐などで呼んで行くのはコストから言って現実的ではありません。
     その解決策としてはプログラムへの統一的な API を提供することができるようすることです。
     awk スクリプトやコンパイラなどを思い浮べるといいと思います。

 \item ユーザへの使い勝手の向上
     ユーザからしても同じ機能なのに別の別の名前であると不便な場合もあります。
     MUA や IRCクライアント からの Web Browser を呼び出す際を想像してもらえば想像し易いと思います。

 \item menu との連携
     Debian 独自なものとして menu がありますが、これと alternatives を連携させることにより
     メニューをいじらずに、違うパッケージを提供することができます。

\end{enumerate}

\subsection{update-alternatives の使い方}
  see man(重要)

  alternatives の制御には /usr/sbin/update-alternatives というコマンドを使います。

\begin{commandline}
  書式一覧
   update-alternatives --install リンク 一般名 パス 優先度  [--slave スレーブリンク スレーブ一般名 スレーブパス]
   update-alternatives --remove 一般名 パス
   update-alternatives --remove-all ディレクトリ
   update-alternatives --all
   update-alternatives --auto 一般名
   update-alternatives --display 一般名
   update-alternatives --list 一般名
   update-alternatives --config 一般名
   update-alternatives --set 一般名 パス
\end{commandline}

\subsection{用語の説明}
\subsubsection{マスター}
\begin{itemize}
 \item 
     マスター  alternatives の対象とするもの
 \item     リンク 一般的なリンクです\\
       例 /usr/bin/awk, /usr/bin/editor, /usr/bin/pager, /usr/bin/x-www-browser
 \item     一般名 一般的な名前です、\\
       例 awk, editor, pager, x-www-browser
 \item     パス それぞれの alternatives のリンク先になります\\
       例 /usr/bin/awk: /usr/bin/nawk, /usr/bin/mawk, /usr/bin/gawk 
          /usr/bin/editor : /bin/ed, /bin/nano, /usr/bin/vim
          /usr/bin/pager : /bin/more, /usr/bin/less, /usr/bin/w3m, /usr/bin/lv
          /usr/bin/x-www-browser : /usr/bin/mozilla, /usr/bin/kazehakase, /usr/bin/mozilla-firefox
 \item     優先度  autoモードで設定される際の選択基準になる整数値です
\end{itemize}
\subsubsection{スレーブ}

\begin{itemize}
 \item 

     スレーブ  alternatives のマスターに付随するファイル、manなど 複数指定可能\\
 \item      スレーブリンク 一般的なリンクに付随するリンク\\
       例 /usr/share/man/man1/awk.1.gz, /usr/share/man/man1/nawk.1.gz, /usr/share/man/man1/editor.1.gz, /usr/share/man/man1/pager.1.gz, /usr/share/man/man1/x-www-browser.1.gz
 \item      スレーブ一般名 一般的な名前に付属する一般名\\
       例 awk.1.gz, editor.1.gz, pager.1.gz, x-www-browser.1.gz
 \item      スレーブパス それぞれの alternatives に付随するリンク先\\
       例 \\ /usr/share/man/man1/awk.1.gz, /usr/share/man/man1/nawk.1.gz : /usr/share/man/man1/mawk.1.gz, /usr/share/man/man1/gawk.1.gz\\
          /usr/share/man/man1/editor.1.gz : /usr/share/man/man1/ed.1.gz, /usr/share/man/man1/nano.1.gz, /usr/share/man/man1/vim.1.gz\\
          /usr/share/man/man1/pager.1.gz : /usr/share/man/man1/more.1.gz, /usr/share/man/man1/less.1.gz, /usr/share/man/man1/w3m.1.gz, /usr/share/man/man1/lv.1.gz\\
          /usr/share/man/man1/x-www-browser.1.gz : /usr/share/man/man1/mozilla.1.gz, /usr/share/man/man1/kazehakase.1.gz, /usr/share/man/man1/mozilla-firefox.1.gz\\
\end{itemize}

\subsubsection{automatic と manual }

        alternatives は優先度によって自動的に選択する automatic モードと
        優先度を無視しユーザが選択する manual モードがあります。 
        デフォルトは automatic モードで --config や --set 、 --all を
        使い変更した場合に manual モードに以降します。
        automatic モードに戻したければ --auto を使います。

\subsection{ユーザとして使う場合}
     ユーザとして主に使用するコマンド

\subsubsection{状態表示(一般ユーザで可)}
\begin{commandline}
     update-alternatives --display 一般名
     update-alternatives --list 一般名
\end{commandline}
\subsubsection{設定変更(rootのみ)}
\begin{commandline}
     update-alternatives --auto 一般名
     update-alternatives --config 一般名
     update-alternatives --set 一般名 パス
\end{commandline}

     の4つです。

\subsubsection{使用例と説明}

\begin{commandline}

     # update-alternatives --display editor             <-- editor という一般名の alternatives のステータスを表示
     editor - status is auto.                           <-- alternatives の モード
      link currently points to /usr/bin/vim             <-- 現在のリンク先
     /bin/ed - priority -100                            <-- マスターリンク - 優先度 (/bin/ed)
      slave editor.1.gz: /usr/share/man/man1/ed.1.gz    <-- スレーブ一般名 : スレーブリンク先 (/bin/ed)
     /bin/nano - priority 40                            <-- マスターリンク - 優先度 (/bin/nano)
      slave editor.1.gz: /usr/share/man/man1/nano.1.gz  <-- スレーブ一般名 : スレーブリンク先 (/bin/nano)
     /usr/bin/vim - priority 120                        <-- マスターリンク - 優先度 (/usr/bin/vim)
      slave editor.1.gz: /usr/share/man/man1/vim.1.gz   <-- スレーブ一般名 : スレーブリンク先 (/usr/bin/vim)
     Current `best' version is /usr/bin/vim.            <-- 優先度の一番高いもの

     # update-alternatives --list editor                <-- editor という一般名の alternatives のマスターリンク先を表示
     /bin/ed                                            <-- マスターリンク先
     /bin/nano                                          <-- マスターリンク先
     /usr/bin/vim                                       <-- マスターリンク先

     # update-alternatives --auto editor                <-- 一般名 editor の alternatives を優先度によって変更
     # update-alternatives --config editor              <-- 一般名 editor の alternatives を選択肢から手動で変更
  
     There are 3 alternatives which provide `editor'.   <-- editor という一般名の選択肢が3つある

        Selection    Alternative
     -----------------------------------------------
          1        /bin/ed                              <-- alternatives
          2        /bin/nano                            <-- alternatives
    *+    3        /usr/bin/vim                         <-- alternatives 現在選択されている

     Press enter to keep the default[*], or type selection number:  <-- ここで番号を選択する
     Using `/usr/bin/vim' to provide `editor'.          <-- editor を /usr/bin/vim で提供する

     # update-alternatives --set editor /usr/bin/vim    <-- 一般名 editor の alternatives を指定したリンクに変更する
     Using `/usr/bin/vim' to provide `editor'.          <-- editor を /usr/bin/vim で提供する
\end{commandline}

\subsection{パッケージメンテナとして使う場合}

     パッケージメンテナな人が主に使うコマンド
     パッケージインストール時に使用(主に postinst)\\
\begin{commandline}
     update-alternatives --install リンク 一般名 パス 優先度  [--slave スレーブリンク スレーブ一般名 スレーブパス]
\end{commandline}

     パッケージ削除時に使用(主に prerm)\\
\begin{commandline}
     update-alternatives --remove 一般名 パス
\end{commandline}

\subsubsection{ 使用例}
     
\begin{commandline}
     update-alternatives --install コマンド

     update-alternatives --install リンク 一般名 パス 優先度  [--slave スレーブリンク スレーブ一般名 スレーブパス]
\end{commandline}

     この書式ですが、slave は省略可能で、使用する場合には複数のスレーブを
     指定することもできます。

     vim の postinst\\
     リンク元のファイルがインストールされてからリンクを貼るので postinst に書く 

\begin{commandline}

     case "$1" in
       abort-upgrade)
         for i in vi view ex editor ; do
           update-alternatives \
             --install /usr/bin/$i $i /usr/bin/vim 120 \
             --slave /usr/share/man/man1/$i.1.gz $i.1.gz /usr/share/man/man1/vim.1.gz

         done
         ;;
       configure)
         for i in vi view ex editor ; do
           update-alternatives \
             --install /usr/bin/$i $i /usr/bin/vim 120 \
             --slave /usr/share/man/man1/$i.1.gz $i.1.gz /usr/share/man/man1/vim.1.gz

         done
         if [ -L /usr/doc/vim ] ; then
           rm /usr/doc/vim
         fi
         ;;
     esac
\end{commandline}     

     vim の prerm\\
     リンク元のファイルがなくなる前に削除するので prerm に書く

\begin{commandline}

     case "$1" in
       remove)
         for i in vi view ex editor ; do
           update-alternatives --remove $i /usr/bin/vim
         done
         ;;
     esac
\end{commandline}

\subsection{マニアックなコマンド}

     update-alternatives --all\\
       全ての選択肢に --config を使い設定を行なう

     update-alternatives --remove-all\\
       ディレクトリにある全ての選択肢を削除する（危険です)

\subsection{一般的なオプション}


\begin{commandline}
     --verbose 冗長メッセージ
     --quiet メッセージを抑制
     --test テスト(未実装)
     --help ヘルプ
\end{commandline}


\subsection{   マニアックなオプション}

     --altdir \\
       リンクを置くディレクトリ\\
       デフォルトは /etc/alternatives/

     --admindir \\
       設定ファイルを置くディレクトリ\\
       デフォルトは /var/lib/dpkg/alternatives/

\subsubsection{ マニアックなオプションの使用例}

     /usr/local/以下にソースからインストールしたものや alternatives を提
     供していないパッケージについて alternatives で管理したい場合などに
     使う 特権ユーザでなくても使えるのも特徴

     使用例:
       /home/foo/eclipse 以下にある /home/foo/eclipse/eclipse を
       alternatives で管理したい。 
\begin{commandline}
       $ mkdir /home/foo/altdir/
       $ mkdir /home/foo/admindir/
       $ mkdir /home/foo/bin/
       $  /usr/sbin/update-alternatives --altdir /home/foo/altdir/ --admindir /home/foo/admindir/ \
              --install /home/foo/bin/eclipse eclipse /home/foo/eclipse/eclipse 100 
\end{commandline}

       で、.bashrc などに /home/foo/bin/ を追加
     IM とか、MUA の管理にも向いているのではないだろうか。



\subsection{update-alternatives ってどうなってるの?}


\subsubsection{   構造}

   /usr/sbin/update-alternatives\\
      alternatives の制御コマンド\\
      perl で実装されており dpkg パッケージに含まれています。

   /var/lib/dpkg/alternatives\\
      alternatives設定ファイルディレクトリ\\
      設定ファイル例

\begin{commandline}
      $ cat /var/lib/dpkg/alternatives/editor  
      alternatives設定ファイルディレクトリ
      /usr/bin/editor                         <-- マスターリンク先
      editor.1.gz                             <-- スレーブ一般名
      /usr/share/man/man1/editor.1.gz         <-- スレーブリンク先
      
      /bin/ed                                 <-- マスターリンク元(/bin/ed)
      -100                                    <-- 優先度(/bin/ed)
      /usr/share/man/man1/ed.1.gz             <-- スレーブリンク元(/bin/ed)
      /bin/nano                               <-- マスターリンク元(/bin/nano)
      40                                      <-- 優先度(/bin/nano)
      /usr/share/man/man1/nano.1.gz           <-- スレーブリンク元(/bin/nano)
      /usr/bin/vim                            <-- マスターリンク元(/usr/bin/vim)
      120                                     <-- 優先度(/usr/bin/vim)
      /usr/share/man/man1/vim.1.gz            <-- スレーブリンク元(/usr/bin/vim)
\end{commandline}

      
   /etc/alternatives/\\
      alternatives のリンクのあるディレクトリです。\\
      リンク例

\begin{commandline}
      $ ls -l /etc/alternatives/editor
      lrwxrwxrwx  1 root root 12 2005-06-01 02:43 /etc/alternatives/editor -> /usr/bin/vi
\end{commandline}

\begin{commandline}
   
   /usr/share/man/man8/update-alternatives.8.gz
   /usr/share/man/de/man8/update-alternatives.8.gz
   /usr/share/man/es/man8/update-alternatives.8.gz
   /usr/share/man/fr/man8/update-alternatives.8.gz
   /usr/share/man/ja/man8/update-alternatives.8.gz
   /usr/share/man/pt_BR/man8/update-alternatives.8.gz
\end{commandline}

      alternativesマニュアルです。
   

\subsubsection{動作}

   update-alternatives --auto\\
\begin{enumerate}
 \item 一般名を取得
 \item      
/var/lib/dpkg/alternatives/一般名 のステータス情報が
        manual の場合は auto に変更
 \item /etc/alternatives/一般名 のリンク先を /var/lib/dpkg/alternatives/一般名 の
        優先度の情報を比較し最大のものへのリンクへと変更
\end{enumerate}

   update-alternatives --config 及び update-alternatives --set コマンド\\
\begin{enumerate}
 \item 一般名とリンク先のパスを取得
 \item /var/lib/dpkg/alternatives/一般名 のステータス情報が
        auto の場合は manual に変更
 \item /etc/alternatives/一般名 のリンク先を指定されたリンク先のパスへと変更
\end{enumerate}

   update-alternatives --install コマンド
\begin{enumerate}
 \item      存在しない場合は指定されたパスのリンクを /etc/alternatives/一般名 へと貼る
 \item      存在しない場合は /var/lib/dpkg/alternatives/一般名 ファイルを指定された情報に基ずいて生成
 \item      /etc/alternatives/一般名 のリンク先を /var/lib/dpkg/alternatives/一般名 の
        優先度の情報を比較し最大のものへのリンクへと変更
\end{enumerate}

   update-alternatives --remove コマンド

\begin{enumerate}
 \item      /var/lib/dpkg/alternatives/一般名 から指定された パスの情報を削除

 \item      /var/lib/dpkg/alternatives/一般名 で alternatives が提供されない場合は
        /var/lib/dpkg/alternatives/一般名 及び /etc/alternatives/一般名
        リンク先を削除

 \item      /var/lib/dpkg/alternatives/一般名 で まだ alternatives が提供されている場合は
        /etc/alternatives/一般名 のリンク先を /var/lib/dpkg/alternatives/一般名 の
        優先度の情報を比較し最大のものへのリンクへと変更
\end{enumerate}


\subsection{update-alternatives の改善案}

   update-alternatives 関連の改善案を自分なりに考えてみた。


\begin{enumerate}
 \item パッケージのファイル情報への反映
     
      現在の alternatives はパッケージをインストールされるまでは
      そのパッケージがどんな alternatives を提案するのか という情報が
      解らない。これだと、インストールしてから、普段使用していた 
      alternatives を使おうとしたら期待してたのと別のものが起動してしまい
      戸惑うことになるし、無駄なハマりの原因となるし、
      もし、そのような情報を集積できれば、機能からパッケージをより
      簡単に検索することができるようになる。

 \item より柔軟に
      
       現在の alternatives では、man を slave にした場合などに
       英語のみしか表示できなくなってしまう、多言語化が進みつつある
       昨今これだけでは不足ということになってしまっているので、
       より柔軟な設定を行なえるのが望ましい。
       ファイルだけではなくコマンドを登録できるようにするなど。

 \item menu 編集インターフェースとの連携

       独自に作成した alternatives などを menu に簡単に登録できるようにする
       などの、応用的な使い方を簡便にするインターフェースの拡充

\end{enumerate}

\subsection{dsys について}

\subsubsection{目的}

   dpkg に含まれる パッケージ依存関係管理を行なう /usr/bin/dpkg 以外の
   /usr/sbin/update-alternatives, /usr/sbin/dpkg-divert, /usr/sbin/dpkg-statoverride
   の3つのコマンドに対応した GUI フロントエンドを提供することを目指している。

\subsubsection{歴史}

   一番初期は update-alternatives のみに対応したもので、ruby 1.6 と ruby-gtk で書いて
   いたが、やっぱ gtk2 だろう！ ということで、ruby-gtk2 に手を出してのんびりしてるうちに
   galternatives など競合で出てきてしまった。
   思い出したら実装したりしているがコーディング能力の低さのために遅々として改善されていない。

\subsubsection{機能}

   /usr/sbin/update-alternatives\\
     一覧の表示、個々のステータスの表示、変更、追加、削除

   /usr/sbin/dpkg-divert\\
     一覧の表示

   /usr/sbin/dpkg-statoverride\\
     一覧の表示、変更、追加、削除

\subsubsection{TODO}

     まずはツールチップを適切に実装すること
     綺麗なコードの書き方や、デザインがわからないのでどうにかする
     せめて man くらい、、
     alternatives と statoverride のインターフェースを意識的に替えていたがそろそろ統一
     treeview で右クリックできるようする
     ショートカットの実装

\subsubsection{将来}

     設定ファイルを作るかも
     アイコン欲しいよね

   協力者の方募集中、、


%% えとーさんの記事はここまで

\dancersection{Inside Debian-Installer}{武藤 健志}
\label{sec:inst}


\subsection{Debian-Installerとは}
\label{sec:whatd-i}

\textbf{Debian-Installer}（以下\textbf{d-i}）は、Debian GNU/Linuxリリース3.1、コードネームSargeから採用された新しいインストーラシステムです。

これまでのインストーラシステム（Woody以前）として採用されていたBoot-Floppiesには、次のような問題がありました。

\begin{itemize}
\item サイズの制限。その名のとおりフロッピーディスクを基盤にしたものなので、システム全体をフロッピー容量の1.2MB〜1.44MBに圧縮して収まるように悪戦苦闘しなければならなかった。
\item 動的ロードの不備。動的にインストーラを拡張する方法がなく、機能向上を行うすべがなかった。
\item ハードウェア認識の不備。ハードウェアを自動認識する機構が入っていなかった。ユーザーはすべて手動で設定しなければならなかった。
\item インストーラのリリースとテスト。インストーラ一式が一体のため、リリースするためには全体をビルドしなければならず、頻繁にリリースすることができなかった。バグへの対処が遅れがちになった。
\end{itemize}

\begin{table}[htbp]
  \begin{tabular}[htbp]{|l|p{10cm}|}\hline
    2000夏ごろ&Joey Hessがデザインを始める\\ \hline
    2000秋ごろ&サンプル実装が公開される\\ \hline
    2002&WoodyがBoot-Floppiesに基づいてリリース。Sargeでd-iになるよう本格的に目標が立てられる\\ \hline
    2003&開発がかなり進展。開発者も増える\\ \hline
    2004&d-i rc2リリース\\ \hline
    2005春&d-i rc3リリース\\ \hline
    2005夏&Sargeリリース\\ \hline
  \end{tabular}
  \caption{d-iの簡単な歴史}
  \label{tab:history}
\end{table}


\subsection{開発体制}
\label{sec:devel}

d-iの開発は、ほかのどのDebianのサブプロジェクトよりも大規模です。

\begin{itemize}
\item リポジトリのコミッタ権限がある人：153人
\item ローカライズコーディネータのコミッタに代行してもらっている人もいるので、実際はもっと多い？
\item オフィシャルなDebian Developerは60人
\item アカウント管理：Alioth（\texttt{\url{http://alioth.debian.org/projects/d-i/}}）、メーリングリスト：debian-boot@lists.debian.org、バグ追跡：installation-reports仮想パッケージ名、IRC：\#debian-boot
\end{itemize}

\begin{table}[htbp]
  \begin{tabular}[htbp]{|l|p{8cm}|}\hline
    プロジェクトリーダ&Joey Hess\\ \hline
    国際化&Christian Perrier、Dennis Stampfer\\ \hline
    フロントエンド&Collin Watson\\ \hline
    パーティションマネージャ&Anton Zinobiev\\ \hline
    ネットワーク設定&Joshua Kwan\\ \hline
    ハードウェア認識&Petter Reinholdtsen\\ \hline
    移植&Martin Michlmayr、Steve Langaek、…\\ \hline
    ドキュメント&Frans Pop\\ \hline
  \end{tabular}
  \caption{主な開発者}
  \label{tab:core}
\end{table}

当初はCVS、現在はSubversionにて協業を行っています。

\begin{itemize}
\item \texttt{\url{http://svn.d-i.alioth.debian.org/svn/d-i/trunk}}（トランク、unstable向け）
\item \texttt{\url{http://svn.d-i.alioth.debian.org/svn/d-i/branches/d-i/sarge}} （フリーズされたSargeのブランチ）
\end{itemize}

\begin{figure}[htbp]
\begin{verbatim}
d-i/
    installer/build/ ――d-iのビルドディレクトリ (make ...を実行するとd-iイメージができる)
    installer/doc/   ――ドキュメント。manual/にXML形式のインストールマニュアルが収録されている
    packages/        ――各udebのソースコード
    scripts/         ――自動化などに使うスクリプト集
\end{verbatim}
  \caption{d-iリポジトリの構造}
  \label{fig:direpo}
\end{figure}

\subsection{d-iの構造}
\label{sec:d-i}

d-iは、これまでに蓄積されてきたさまざまなLinux/Debianのフレームワークを応用しています。

\begin{description}
\item[udeb] 通常のdebから実行に不必要なものを削除してスリムにしたもの
\item[cdebconf] Cで記述されたdebconfフレームワーク (いくらか拡張)
\item[devfs] ハードディスクなどのデバイスファイルの取り扱いを容易にする
\item[フレームバッファと国際化端末] 国際化端末を実行する (bogl-termおよびjfbterm)
\item[discoverとhotplug] ハードウェアの自動認識
\end{description}


\subsubsection{udeb}
\label{sec:udeb}

udebは、通常のdebとほぼ同じですが、ドキュメントなどの実行に直接関係ないものを削ぎ落としたバイナリパッケージです。

\begin{itemize}
\item インストーラだけで利用する特別なメタ情報として、XB-Installer-Menu-Itemが含まれているものがある。これは、インストーラのメニューに表示する際の順番を示す
\item DependsやProvidesで依存関係を処理したり、debconf用の国際化済みtemplatesによってローカライズされたメニューを表示できる (メニューに表示されるのはパッケージのshort descriptionだが、特別にdebconf templatesの中でこのローカライズも設定している)
\item インストーラのメニューを選択することは、postinstを実行するということになる
\item 環境によってロードする順序を変えることができる 。たとえば、ネットワークインストールなら最初にネットワークの設定を済ませてからIDEコントローラを検出、CDインストールなら最初にIDEコントローラの設定を済ませてからネットワークを検出
\end{itemize}

\begin{table}[htbp]
  \begin{tabular}[htbp]{|l|p{7cm}|}\hline
    anna&擬似APT\\ \hline
    udpkg&マイクロdpkg\\ \hline
    languagechooser、countrychooser&言語・国情報の設定\\ \hline
    hw-detect、ethdetect、hw-detect-full&ハードウェアの検出\\ \hline
    preseed&事前構成のロード\\ \hline
    netcfg&ネットワーク設定\\ \hline
    network-console&SSHサーバーを起動してリモートインストールを実行する\\ \hline
    cdrom-retriever、floppy-retriever、net-retriever&udebパッケージをダウンロードする\\ \hline
    partman&パーティションマネージャ\\ \hline
    base-installer&ベースシステムをdebootsrapで展開し、CPUに合ったカーネルをインストールする\\ \hline
    os-prober&インストールされているほかのOSの検出\\ \hline
    prebaseconfig&CDの取り出し、現在の状態の保存など\\ \hline
  \end{tabular}
  \caption{主なudeb}
  \label{tab:udeb}
\end{table}

リポジトリのpackages/だけでなく、libcやdiscover、console-toolsのように、外部のdebパッケージがudebを提供することもあります。

d-iはudebの集合体で、最初にロードされるインストーラもudebを基に構成されています。d-iのmake時にどのudebから構成しておくかによって、CD向け、ネットワーク向け、USBメモリ向けといったイメージを作成できます。


\subsubsection{cdebconf}
\label{sec:cdebconf}

cdebconfは、Perlで記述されていたdebconfを、Cで実装し直したものです。

\begin{itemize}
\item Perlの膨大なアーカイブの必要なく動作
\item 進捗バーなどの機能拡張
\end{itemize}

今後は普通の環境も、debconfは捨てて(orphan済み)、cdebconfに移行しようという計画になっています。

d-iではcdebconfが大きな役割を果たしています。

\begin{itemize}
\item debconfインターフェイスにのっとったフロントエンドを切り替えることができる (デフォルトはdialog)
\item templatesからインストールの質問を表示し、答えをデータベースに格納する
\item 質問に優先度(critical、high、medium、low)を設定して、質問の詳細と数を変えられる (デフォルトはhigh)
\item データベースに最初から値を入れるようにして、全自動インストールができる (preseed)
\end{itemize}


\subsubsection{devfs}
\label{sec:devfs}

Linuxカーネル2.4で導入されたdevfsでは、「動的に」「利用可能な」デバイスファイルが作成されるので、ハードディスクのパーティションやフレームバッファの存在を検出するのに便利です。

ただ、devfsはもうobsoleteなので、今後はudevに書き換えていくことになるでしょう。


\subsubsection{フレームバッファと国際化端末}
\label{sec:i18n}

Linuxのデフォルトの端末画面はラテン文字しか表示できないので、国際化されたインストーラを実行するためにはフレームバッファドライバをロードしたあと、その上で国際化端末を実行する、という手順が必要になります。x86の場合フレームバッファドライバは、vesafb→vga16fbと試行しています。d-iの国際化端末としては、次の2つを採用しています。

\begin{itemize}
\item bogl-bterm：UTF-8対応の端末。
\item jfbterm：再起動後、CJK圏でのみ使われる端末(表示はEUC-JP)。なぜbogl-btermで続けなかったかというと…
\end{itemize}


\subsubsection{discoverとhotplug}
\label{sec:discover}

最初のハードウェア認識にはdiscoverが使われます。再起動後、標準で使われるのはhotplugです。

\begin{itemize}
\item discover：能動的にハードウェアを認識する。データベースにはPCI ID値などを列挙したdiscover1-dataパッケージの情報を使う
\item hotplug：ほぼ受動的にハードウェアを認識する。カーネルから送られた情報を使う
\end{itemize}


\subsection{今後のd-i}
\label{sec:di}

Etchに向けてのTODOはいろいろあります。

\begin{itemize}
\item より良いハードウェア認識：discover2への移行。既存のdiscoverのデータベースの更新(volatile経由？)。
\item Woodyのときに使えたような「レスキューモード」：rescue.udebが用意されました
\item Gtk+によるGUIフロントエンド：実装が進んでいます。ただ、udebの持っている情報をもっと増やすか、別のGUI専用のファイルを用意しないと、GUIのメリットを活かした画面にはならないと思います
\item 複数言語の選択：だいぶ実装されてきてはいるようです
\item フロッピーの取り扱い：今回は技術上、カーネル2.4のみの対応としています(カーネル2.6はフロッピーに入らない)。カーネルを分割してロードできないか？
\end{itemize}

\newpage
%\vfill{}
%\hfill{}
%\includegraphics[width=7cm]{image200502/openlogo-nd.eps}
%\hfill{}
%\vfill{}
%\newpage

\dancersection{個人提案課題}{}
\hfill{}{\large 名前} \underline{\hspace{6cm}}

下記の空欄を埋めてください:

{\LARGE 
Debian installer の (\hspace{5cm})\\
に注目し、今後のネタは(\hspace{6cm})\\
します。
}
\\\\
企画案の図：

\newpage

%\vfill{}
%\hfill{}
%\includegraphics[width=7cm]{image200502/openlogo-nd.eps}
%\hfill{}
%\vfill{}
%\newpage

\dancersection{Keysigning Party}{上川 純一}


事前に必要なもの
\begin{itemize}
 \item 自分の鍵のfingerprintを書いた紙({\tt gpg --fingerprint XXXX}の出力。)
 \item 写真つきの公的機関の発行する身分証明書、fingerprintに書いてある名前が自分の
ものであると証明するもの
\end{itemize}

キーサインで確認する内容
\begin{itemize}
 \item 相手が主張している名前の人物であることを信頼できる身分証明書で証明し
       ているか\footnote{いままで見た事の無い種類の身分証明書を見せられ
       てもその身分証明書の妥当性は判断しにくいため、学生証明書やなんと
       か技術者の証明書の利用範囲は制限される。
       運転免許証明書やパスポートが妥当と上川は判断している}。
 \item 相手がfingerprintを自分のものだと主張しているか
 \item 相手のfingerprintに書いてあるメールアドレスにメールをおくって、
       その暗号鍵にて復号化することができるか
\end{itemize}

手順としては
\begin{itemize}
 \item 相手の証明書を見て、相手だと確認
 \item fingerprintの書いてある紙をうけとり、これが自分のfingerprintだと
       いうことを説明してもらう
 \item (後日) gpg署名をしたあと、鍵のメールアドレスに対して暗号化して送付、相手
       が復号化してキーサーバにアップロードする
       ({\tt gpg --sign-key XXXXX, gpg --export --armor XXXX })
\end{itemize}


\dancersection{次回}{}

次回は7月2日土曜日の夜を予定しています。
内容は本日決定予定です。

参加者募集はまた後程。

\end{document}